\chapter[Voorbeelden van codefragmenten]{Voorbeelden van code-\\fragmenten}
Hier is een stukje C code:
\begin{lstfragment}[style=cstyle]
int main() {
    printf("Hello LaTeX");
    return 0;
}
\end{lstfragment}
\begin{clstShortInline}
De functie |main|\footnote{%
    In een C-programma wordt de functie \ccode{main} als eerste aangeroepen.
} returned de waarde |0| met een |return|-statement.
 Het woord |class| is geen keyword in C.
\end{clstShortInline}

Hier is een stukje \Cpp{} code:
\begin{lstfragment}[style=cstyle]
// voorbeeld van C++ code
int main() {
    std::cout << "Hello LaTeX" << std::endl;
    return 0;
}
\end{lstfragment}
\begin{cpplstShortInline}
De functie |main|\footnote{%
    Ook in een \Cpp{}-programma wordt de functie \cppcode{main} als eerste aangeroepen.
} returned de waarde |0| met een |return|-statement.
Het woord |class| is een keyword in \Cpp{}.
\end{cpplstShortInline}

Hier is een stukje MATLAB-code:
\begin{lstfragment}[style=mstyle]
% voorbeeld van MATLAB code
b = 0.25 * sinc(0.25 * [-10:10]);
[H,f] = freqz(b, 1, 512, 8000);
HdB = 20 * log10(abs(H));
plot(f, HdB, 'Color', [204/255 0 51/255]);
grid;
ylabel('Magnitude (dB)');
ylim([-45 5]);
xlabel('Frequency (Hz)');
\end{lstfragment}

\clearpage
\chapter{Voorbeelden van floats}

\Cref{lst:fdacoefs} is een door MATLAB gegenereerde headerfile.
De listing is floating  en kan dus ergens anders geplaatst worden dan meteen onder deze tekst.

\floatlstinput[style=cstyle,style=linenumbers]{fdacoefs.h}{fdacoefs}{16-bit signed integer coefficients generated by MATLAB.}

\Cref{fig:binary} laat een mooi plaatje zien als een floating figuur en kan dus ergens anders geplaatst worden dan meteen onder deze tekst.

\figuur{width=.5\textwidth}{binary}{Nullen en enen --- het mooiste wat er is.}

\Cref{tab:simpel} is een voorbeeld van een simpele floating tabel en kan dus ergens anders geplaatst worden dan meteen onder deze tekst.

\tabel{Een eenvoudige vertaaltabel.}{simpel}{cc}{
    Engels & Nederlands
}{
    biscuit & kaakje \\
    cake & cake \\
    holiday & vakantie
} 

Deze tekst komt (misschien) onder de tabel zodat de afstand tussen de tabel en de tekst zichtbaar wordt.

\Cref{lst:ledcopy} is een eenvoudig AVR-assembler programma.
De listing is floating en kan dus ergens anders geplaatst worden dan meteen onder deze tekst.

\floatlstinput[style=astyle,style=linenumbers,linerange={1-2,8-17}]{ledcopy.asm}{ledcopy}{Kopiëren van inputs naar outputs in AVR-assembler.}

\clearpage
\chapter{Voorbeelden van Here floats}
Here floats zijn helemaal geen floats want ze worden op de plaats gezet waar ze zijn gedefinieerd.
Dit kan ongewenste witruimte aan het einde van een pagina tot gevolg hebben.
Het verdient dan ook de voorkeur om hier alleen in uiterste nood gebruik van te maken.
\LaTeX{} weet vaak beter waar figuren, tabellen en listings geplaatst moeten worden om een mooie bladspiegel te creëren.
\Cref{lst:fdacoefs2} is een door MATLAB gegenereerde headerfile.
Deze listing komt gegarandeerd meteen na deze tekst.

\lstinput[style=cstyle,style=linenumbers]{fdacoefs.h}{fdacoefs2}{16-bit signed integer coefficients generated by MATLAB.}

\Cref{fig:fractal} laat een mooi figuur zien en komt gegarandeerd meteen na deze tekst.

\figuur[H]{width=.5\textwidth}{fractal}{Fractals --- ook mooi.}

\Cref{tab:simpel2} is een voorbeeld van een simpele tabel en komt gegarandeerd meteen na deze tekst.

\tabel[H]{Een eenvoudige vertaaltabel.}{simpel2}{cc}{
    Engels & Nederlands
}{
    biscuit & kaakje \\
    cake & cake \\
    holiday & vakantie
} 

Deze tekst komt onder de tabel zodat de afstand tussen de tabel en de tekst zichtbaar wordt.

